\section{Method}

\subsection{Overview}

JambaInfo is a two-branch hybrid block:

\begin{itemize}
\item \textbf{Branch 1: Differential attention} (noise-reduced).
\item \textbf{Branch 2: Concept-bottlenecked SSM} (IB-motivated compression).

\end{itemize}
The branches are fused by a learned scalar gate.

\textbf{Figure 4.1} shows the pipeline.

\begin{lstlisting}
    Input X $\in$ R^{B$\times$L$\times$d}
              |
       +------+------+
                      
   +-------+    +---------------+
   | Diff  |    |  Bottleneck   |
   | Attn  |    |  W\_down $\to$     |
   |       |    |  GELU $\to$       |
   |       |    |  W\_up         |
   +---+---+    +-------+-------+
       |                 
       |            +-------+
       |            |  SSM  |
       |            +---+---+
       | $\alpha$=$\sigma$(w)         |
       |  +-------------+
           
   (1-$\alpha$) A + $\alpha$ M
       |
        
     out\_proj
       |
        
    Output Y
\end{lstlisting}

\subsection{Differential Attention Branch}

Same as HymbaTransX:

\[
Q\_1 = X W\_{Q\_1}, \quad K\_1 = X W\_{K\_1}, \quad
Q\_2 = X W\_{Q\_2}, \quad K\_2 = X W\_{K\_2}
\tag{4.1}
\]

\[
A\_1 = \text{softmax}(Q\_1 K\_1^\top / \sqrt{d\_h}), \quad
A\_2 = \text{softmax}(Q\_2 K\_2^\top / \sqrt{d\_h})
\tag{4.2}
\]

\[
A\_\text{diff} = \text{ReLU}(A\_1 - 0.2 A\_2), \quad
\tilde{A}\_\text{diff} = \frac{A\_\text{diff}}{\sum A\_\text{diff} + \epsilon}
\tag{4.3}
\]

\[
A\_\text{out} = \tilde{A}\_\text{diff} V, \quad V = X W\_V
\tag{4.4}
\]

\subsection{Concept-Bottlenecked SSM Branch}

\textbf{Step 1 -- Bottleneck projection:}

\[
z = W\_\text{up} \cdot \text{GELU}(W\_\text{down} X)
\tag{4.5}
\]

with $W_\text{down} \in \mathbb{R}^{b \times d}$, $W_\text{up} \in \mathbb{R}^{d \times b}$.

\textbf{Step 2 -- SSM on concept representation:}

For decoders:
\[
M\_\text{concept} = \text{SSM}(z)
\tag{4.6}
\]

For encoders (bidirectional):
\[
M\_\text{fwd} = \text{SSM}(z), \quad M\_\text{bwd} = \text{flip}(\text{SSM}(\text{flip}(z)))
\tag{4.7}
\]

\[
M\_\text{concept} = 0.5(M\_\text{fwd} + M\_\text{bwd})
\tag{4.8}
\]

\subsection{Why a Concept Bottleneck?}

\subsubsection{Information-Theoretic Motivation}

The SSM state $h_t \in \mathbb{R}^d$ has bounded capacity. Feeding it a 
$d$-dimensional input with high redundancy wastes capacity on noise. The 
bottleneck reduces the input to $b = 256$ dimensions of "concepts", 
concentrating information.

\subsubsection{Parameter Reduction}

Without bottleneck: SSM has $d$-dimensional input and output.
With bottleneck: SSM still $d \times d$, but the input to SSM has fewer 
unique directions.

\subsubsection{Empirical Observation}

\textbf{Figure 4.2} (deferred to Section 8) shows that the SSM's recurrent 
state has effective rank closer to $b$ than $d$ -- confirming the 
bottleneck dimension is appropriate.

\subsection{Fusion}

\[
\alpha = \sigma(w), \quad w\_0 = -2
\tag{4.9}
\]

\[
\text{combined} = (1 - \alpha) \odot A\_\text{out} + \alpha \odot M\_\text{concept}
\tag{4.10}
\]

\textbf{Note the swap.} HymbaTransX uses $\alpha A + (1-\alpha) M$; we use 
$(1-\alpha) A + \alpha M$. Combined with $\alpha_0 = 0.12$, this favors 
attention initially, letting the model learn how much SSM it needs.

\subsection{Dual Operational Modes}

JambaInfo supports two modes:

\begin{table}[H]
\centering
\small
\begin{tabular}{llll}
\toprule
\textbf{Mode} & \textbf{Description} & \textbf{$\alpha$ init} & \textbf{Use case} \\
\midrule
\texttt{retrieval} & Default. Bias toward attention. & -2 & Classification, QA \\
\texttt{compression} & Bias toward SSM. & +2 & Long-context, generation \\
\bottomrule
\end{tabular}
\end{table}

The \texttt{mode} parameter is passed at construction; the difference is the 
initial value of $w$ and a slight rescaling of the SSM output.

\subsection{Sparsity Parameter}

The \texttt{sparsity} parameter (default 0.5) controls an optional sparsity 
penalty on the differential attention. When active, it encourages the 
$A_\text{diff}$ map to be sparse (few nonzero entries).

In the evaluated architecture, \texttt{sparsity} establishes the coefficient for future regularized $L_1$ differential attention pruning passes.

\subsection{Algorithm}

\textbf{Algorithm 1:} JambaInfo forward pass
\begin{lstlisting}
Input:  X $\in$ R^{B$\times$L$\times$d}, attention\_mask
        position\_embeddings
Output: Y $\in$ R^{B$\times$L$\times$d}

 1: Q1, K1 $\leftarrow$ W\_q1(X), W\_k1(X)
 2: Q2, K2 $\leftarrow$ W\_q2(X), W\_k2(X)
 3: V $\leftarrow$ W\_v(X)
 4: if pos\_emb: apply RoPE
 5: if GQA: repeat\_kv
 6: A1 $\leftarrow$ softmax(Q1 K1^T / $\sqrt{}$d\_h)
 7: A2 $\leftarrow$ softmax(Q2 K2^T / $\sqrt{}$d\_h)
 8: A\_diff $\leftarrow$ ReLU(A1 - 0.2 A2)
 9: A\_out $\leftarrow$ normalize(A\_diff) @ V
10: z $\leftarrow$ W\_up(GELU(W\_down(X)))         Concept
11: if decoder: M $\leftarrow$ SSM(z)
12: else: M $\leftarrow$ 0.5(SSM(z) + flip(SSM(flip(z))))
13: $\alpha$ $\leftarrow$ sigmoid(w)
14: out $\leftarrow$ (1-$\alpha$) A\_out + $\alpha$ M
15: Y $\leftarrow$ out\_proj(out)
16: return Y
\end{lstlisting}

\subsection{Complexity Analysis}

Per block:

\begin{table}[H]
\centering
\small
\begin{tabular}{lll}
\toprule
\textbf{Component} & \textbf{Time} & \textbf{Params} \\
\midrule
Q/K/V & $O(L d^2)$ & $3 d^2$ \\
W\_q2, W\_k2 & $O(L d^2)$ & $2 d^2$ \\
Diff attention & $O(L^2 d)$ & 0 \\
Bottleneck & $O(L d b)$ & $2 b d$ \\
SSM & $O(L d)$ & $2 d^2$ \\
out\_proj & $O(L d^2)$ & $d^2$ \\
\textbf{Total} & $O(L^2 d)$ & $8 d^2 + 2bd$ \\
\bottomrule
\end{tabular}
\end{table}

Compared to baseline ($4 d^2$), we have ~2$\times$ more parameters.

\subsection{Parameter Analysis}

For $d = 768$, $b = 256$:

\begin{table}[H]
\centering
\small
\begin{tabular}{lll}
\toprule
\textbf{Component} & \textbf{Baseline} & \textbf{JambaInfo} \\
\midrule
Q, K, V, O & 4$\times$589,824 = 2.36M & -- \\
W\_q1, W\_k1 & -- & 1.18M \\
W\_q2, W\_k2 & -- & 1.18M \\
W\_v & -- & 0.59M \\
Bottleneck & -- & 0.39M \\
SSM & -- & 0.59M \\
W\_o & -- & 0.59M \\
\textbf{Total} & \textbf{2.36M} & \textbf{4.52M} \\
\textbf{Increase} & -- & \textbf{+92\%} \\
\bottomrule
\end{tabular}
\end{table}

\subsection{Training}

Trainable:
\begin{itemize}
\item LoRA adapters on all attention projections.
\item $W_{Q_2}, W_{K_2}$ (Kaiming init).
\item Bottleneck $W_\text{down}, W_\text{up}$.
\item SSM parameters.
\item $\alpha$ (scalar).

\end{itemize}
Frozen:
\begin{itemize}
\item Base Q/K/V/O weights.

\end{itemize}
\textbf{Trainable count:} ~0.8M per layer.
\begin{lstlisting}

---
\end{lstlisting}